% Copyright (c) 2026 Eugene Joseph M. Ragasa.
% SPDX-License-Identifier: Apache-2.0
% Author: Eugene Joseph M. Ragasa
% Formal proof completed: 2026-10-19
\documentclass[11pt]{article}
\usepackage[T1]{fontenc}
\usepackage{amsmath,amssymb,booktabs,tabularx}
\usepackage[margin=1in]{geometry}
\title{Proof Dossier: Minimum Feasible Cost Is Attained}
\author{Eugene Joseph M. Ragasa}
\date{Formal proof completed: 2026-10-19}
\begin{document}\maketitle
\section{Identity and status}
\begin{tabularx}{\textwidth}{@{}lX@{}}\toprule
Statement ID & \texttt{BG-DECISION-THM-011} \\
Status & proved; no admissions \\
Lean declaration & \texttt{ProjectKoios.BandGauge.Decision.exists\_minimumFeasibleCost} \\
Owning module & \texttt{ProjectKoios/BandGauge/Decision/FinitePareto.lean} \\
Axioms & \texttt{propext}, \texttt{Classical.choice}, \texttt{Quot.sound} \\
\bottomrule\end{tabularx}
\section{Statement}
For a finite decision type, a linearly ordered cost, and a nonempty feasible class, some
feasible decision has cost no larger than every other feasible decision.
\section{Argument}
The feasible subset is finite and nonempty. Choose a cost-minimal member. In a linear
order, if its cost were not at most another feasible cost, the other cost would be
strictly smaller, contradicting minimality.
\section{Boundary and provenance}
The theorem assigns no finite value to an empty feasible class and is exhaustive only for
the declared finite type. Lean checked the exact statement; the dossier and proof were
AI-assisted.
\end{document}
